\documentclass{article}
\usepackage{amsmath,amssymb,amsthm}
\usepackage{algorithm}
\usepackage{algorithmic}
\usepackage{enumitem}
\usepackage{hyperref}
\newtheorem{lemma}{Lemma}
\newtheorem{theorem}{Theorem}
\begin{document}

%=====================================================================
\section*{Algorithm Name}
\textbf{Interpolated Snapshot Averaging (ISA)}  

ISA runs a base stochastic optimizer (e.g. SGD) and stores the model
parameters every $k$ iterations.  At test time an \emph{interpolated}
model is obtained by linearly blending the two most recent stored
snapshots.  The method is designed for non‑convex smooth objectives and
inherits the stability of weight averaging while allowing a cheap
interpolation to any intermediate point.

%=====================================================================
\section*{Mathematical Formulation}
Consider the stochastic optimisation problem  

\[
\min_{w\in\mathbb{R}^{d}} \; f(w):=\mathbb{E}_{\xi\sim\mathcal{D}}[F(w;\xi)] ,
\]

where $F(\cdot;\xi)$ is differentiable and possibly non‑convex.
Let $\{g_t\}_{t\ge 1}$ be stochastic gradients
\[
g_t:=\nabla F(w_t;\xi_t),\qquad \xi_t\stackrel{i.i.d.}{\sim}\mathcal{D}.
\]

\medskip
\noindent\textbf{Base update (SGD).}  
\[
\boxed{w_{t+1}=w_t-\eta_t g_t},\qquad t=0,1,\dots,T-1,
\tag{1}
\]
where $\eta_t>0$ is a stepsize (possibly time‑varying).

\medskip
\noindent\textbf{Snapshot rule.}  
Define the set of snapshot indices  

\[
\mathcal{S}:=\{s\in\mathbb{N}\;|\; s=k\ell,\ \ell=0,1,\dots,\lfloor T/k\rfloor\}.
\]

For every $s\in\mathcal{S}$ store  
\[
\boxed{w^{(s)}:=w_{s}} .
\tag{2}
\]

\medskip
\noindent\textbf{Interpolation for prediction.}  
Let $s_{-}:=\max\{s\in\mathcal{S}\;|\;s\le \tau\}$ and
$s_{+}:=\min\{s\in\mathcal{S}\;|\;s\ge \tau\}$ for a query time
$\tau\in[0,T]$.
Define the interpolation coefficient  

\[
\lambda:=\frac{\tau-s_{-}}{s_{+}-s_{-}}\in[0,1].
\]

The interpolated weight is  

\[
\boxed{\tilde w(\tau)= (1-\lambda) w^{(s_{-})}+ \lambda\, w^{(s_{+})}} .
\tag{3}
\]

In particular, for the final output we set $\tau=T$ and use the two
latest snapshots $w^{(s_{-})},w^{(s_{+})}$.

%=====================================================================
\section*{Pseudocode}
\begin{algorithm}[H]
\caption{Interpolated Snapshot Averaging (ISA)}
\begin{algorithmic}[1]
\STATE \textbf{Input:} initial point $w_0$, stepsize schedule $\{\eta_t\}$,
snapshot period $k$, total iterations $T$.
\STATE Initialize $w \leftarrow w_0$.
\FOR{$t=0$ \TO $T-1$}
    \STATE Sample $\xi_t\sim\mathcal{D}$ and compute $g_t=\nabla F(w;\xi_t)$.
    \STATE $w \leftarrow w - \eta_t g_t$ \hfill\COMMENT{SGD update (1)}
    \IF{$t$ mod $k =0$}
        \STATE Store snapshot $w^{(t)}\leftarrow w$ \hfill\COMMENT{Rule (2)}
    \ENDIF
\ENDFOR
\STATE Let $s_{-}=\max\{s\in\mathcal{S}\mid s\le T\}$,
      $s_{+}=\min\{s\in\mathcal{S}\mid s\ge T\}$ (if $T$ itself is a snapshot,
      $s_{-}=s_{+}=T$).
\STATE $\lambda\gets \dfrac{T-s_{-}}{s_{+}-s_{-}}$ (set $\lambda=0$ if $s_{-}=s_{+}$).
\STATE \textbf{Output: } $\displaystyle \tilde w = (1-\lambda)w^{(s_{-})}+ \lambda w^{(s_{+})}$ \hfill\COMMENT{Interpolation (3)}
\end{algorithmic}
\end{algorithm}

%=====================================================================
\section*{Dry Run Example}
 minimise $f(w)=(w-3)^2$ (one‑dimensional).  
Gradient: $\nabla f(w)=2(w-3)$.  
Set $w_0=0$, constant stepsize $\eta=0.1$, snapshot period $k=2$, run $T=3$.

\begin{enumerate}[label=\arabic*.]
\item $t=0$: $g_0=2(0-3)=-6$  
      $w_1 = 0 -0.1(-6)=0.6$  
      $f(w_1)=(0.6-3)^2=5.76$  
      ($t=0$ is a snapshot because $0\!\!\mod 2=0$) → store $w^{(0)}=0.6$.
\item $t=1$: $g_1=2(0.6-3)=-4.8$  
      $w_2 = 0.6 -0.1(-4.8)=1.08$  
      $f(w_2)=(1.08-3)^2=3.6864$  
      ($t=1$ not a snapshot).
\item $t=2$: $g_2=2(1.08-3)=-3.84$  
      $w_3 = 1.08 -0.1(-3.84)=1.464$  
      $f(w_3)=(1.464-3)^2=2.3665$  
      $t=2$ is a snapshot → store $w^{(2)}=1.464$.
\end{enumerate}

Final query time $\tau=3$.  
$s_{-}=2$, $s_{+}=2$ (no later snapshot), thus $\lambda=0$ and  

\[
\tilde w = w^{(2)} = 1.464,\qquad f(\tilde w)=2.3665 .
\]

The interpolated model lies between the two stored points and has
already moved considerably towards the optimum $w^\star=3$.

%=====================================================================
\section*{Assumptions}
\begin{enumerate}[label=(A\arabic*)]
\item \textbf{Smoothness.} $f$ is $L$‑smooth:
\[
\|\nabla f(x)-\nabla f(y)\|\le L\|x-y\|,\qquad\forall x,y\in\mathbb{R}^d .
\tag{A1}
\]
\item \textbf{Unbiased stochastic gradients.}
\[
\mathbb{E}[g_t\mid w_t]=\nabla f(w_t),\qquad\forall t .
\tag{A2}
\]
\item \textbf{Bounded variance.}
\[
\mathbb{E}\big[\|g_t-\nabla f(w_t)\|^2\mid w_t\big]\le\sigma^2 ,
\qquad\forall t .
\tag{A3}
\]
\item \textbf{Bounded second moment of gradients.}
\[
\mathbb{E}\big[\|g_t\|^2\mid w_t\big]\le G^2 ,\qquad\forall t .
\tag{A4}
\]
\item \textbf{Stepsize.} Constant stepsize $\eta_t\equiv\eta\le\frac{1}{L}$.
\tag{A5}
\item \textbf{Snapshot period.} $k\ge 1$ is a fixed integer independent of $T$.
\tag{A6}
\end{enumerate}

%=====================================================================
\section*{Main Convergence Theorem}
\begin{theorem}[Convergence of ISA]\label{thm:isa}
Under assumptions (A1)–(A6) and after $T$ SGD iterations with
constant stepsize $\eta\le \frac{1}{L}$,
let $\tilde w$ be the interpolated model defined in (3) using the two
latest snapshots.  Then  

\[
\boxed{\;
\mathbb{E}\big[\|\nabla f(\tilde w)\|^2\big]\;
\le\;
\frac{2\big(f(w_0)-f^{\star}\big)}{\eta T}
\;+\;
L\eta G^2
\;+\;
\frac{4L^2\eta^2\sigma^2}{k}
\;}
\tag{4}
\]

where $f^{\star}=\inf_{w}f(w)$.  Consequently, choosing
$\eta=\Theta\!\big(\tfrac{1}{\sqrt{T}}\big)$ yields  

\[
\mathbb{E}\big[\|\nabla f(\tilde w)\|^2\big]=
\mathcal{O}\!\big(T^{-1/2}\big).
\]
\end{theorem}

%=====================================================================
\section*{Theoretical Properties}
\begin{itemize}
\item \textbf{Stability.} The interpolation (3) is a convex combination of
two iterates that are $k$ steps apart; smoothness guarantees that the
gradient does not change abruptly, so the interpolated point inherits the
stability of each snapshot.

\item \textbf{Hyperparameter sensitivity.}
    \begin{itemize}
    \item Larger $k$ reduces the number of stored models but increases the
          third term in (4); the bound grows as $\mathcal{O}(1/k)$.
    \item Stepsize $\eta$ appears in both a decreasing ($1/(\eta T)$) and an
          increasing ($L\eta$) term, leading to the optimal choice
          $\eta\asymp T^{-1/2}$.
    \end{itemize}
\item \textbf{Comparison to baselines.}
    \begin{itemize}
    \item Standard SGD without averaging yields the same first two terms of
          (4) but lacks the $1/k$ reduction term that arises from
          averaging across snapshots.
    \item Classical Stochastic Weight Averaging (SWA) averages *all*
          iterates; ISA uses only two snapshots, giving comparable
          guarantees with far less memory and the additional flexibility
          of interpolation.
    \end{itemize}
\end{itemize}

%=====================================================================
\section*{Discussion}
The theorem shows that ISA attains the same $\mathcal{O}(T^{-1/2})$
convergence rate for the squared gradient norm as plain SGD, while
the interpolation term provides a controllable bias that vanishes as the
snapshot period $k$ grows.  In practice, modest values such as $k\in[5,20]$
lead to smoother test‑time predictions without sacrificing the theoretical
guarantee.

%=====================================================================
\section*{Preliminaries (Definitions and Lemmas)}
\begin{lemma}[Smoothness inequality]\label{lem:smooth}
If $f$ is $L$‑smooth, then for any $x,y$,
\[
f(y)\le f(x)+\langle\nabla f(x),y-x\rangle+\frac{L}{2}\|y-x\|^{2}.
\tag{5}
\]
\end{lemma}
\begin{lemma}[One‑step descent for SGD]\label{lem:descent}
Under (A1)–(A5) and update (1),
\[
\mathbb{E}\big[f(w_{t+1})\mid w_t\big]\le
f(w_t)-\frac{\eta}{2}\|\nabla f(w_t)\|^{2}
+\frac{L\eta^{2}}{2}\mathbb{E}\big[\|g_t\|^{2}\mid w_t\big].
\tag{6}
\]
\end{lemma}
\begin{proof}
Apply Lemma~\ref{lem:smooth} with $x=w_t$, $y=w_{t+1}=w_t-\eta g_t$:
\[
f(w_{t+1})\le f(w_t)-\eta\langle\nabla f(w_t),g_t\rangle
+\frac{L\eta^{2}}{2}\|g_t\|^{2}.
\]
Take conditional expectation and use (A2):
\[
\mathbb{E}[\,\langle\nabla f(w_t),g_t\rangle\mid w_t]
= \|\nabla f(w_t)\|^{2}.
\]
Hence (6). \qedhere
\end{proof}

%=====================================================================
\section*{Main Proof (Step-by-Step Derivation)}
We prove Theorem~\ref{thm:isa}.  All expectations are taken over the
randomness of the stochastic gradients and the snapshot indices.

\begin{enumerate}[label=[Step \arabic*]]
\item \textbf{Summation of one‑step descent.}  
    Sum inequality (6) from $t=0$ to $T-1$ and take total expectation:
    \[
    \sum_{t=0}^{T-1}\frac{\eta}{2}\mathbb{E}\big[\|\nabla f(w_t)\|^{2}\big]
    \le f(w_0)-\mathbb{E}[f(w_T)]
    +\frac{L\eta^{2}}{2}\sum_{t=0}^{T-1}\mathbb{E}\big[\|g_t\|^{2}\big].
    \tag{7}
    \]
\item \textbf{Bounding gradient second moments.}  
    From (A4), $\mathbb{E}[\|g_t\|^{2}]\le G^{2}$, thus
    \[
    \frac{L\eta^{2}}{2}\sum_{t=0}^{T-1}\mathbb{E}[\|g_t\|^{2}]
    \le \frac{L\eta^{2}G^{2}T}{2}.
    \tag{8}
    \]
\item \textbf{Relating snapshots to iterates.}  
    Let $S:=\lfloor T/k\rfloor$ be the number of stored snapshots.
    Define the average of snapshots
    \[
    \bar w_{\text{snap}}:=\frac{1}{S+1}\sum_{\ell=0}^{S} w^{(\ell k)} .
    \tag{9}
    \]
    By convexity of the squared norm,
    \[
    \big\|\nabla f(\bar w_{\text{snap}})\big\|^{2}
    \le \frac{1}{S+1}\sum_{\ell=0}^{S}\|\nabla f(w^{(\ell k)})\|^{2}.
    \tag{10}
    \]
\item \textbf{Bounding the average of snapshot gradients.}  
    Since each snapshot $w^{(\ell k)}$ coincides with an iterate $w_t$,
    we may bound the right‑hand side of (10) by the average of all
    iterates:
    \[
    \frac{1}{S+1}\sum_{\ell=0}^{S}\|\nabla f(w^{(\ell k)})\|^{2}
    \le \frac{k}{T}\sum_{t=0}^{T-1}\|\nabla f(w_t)\|^{2}.
    \tag{11}
    \]
    (Every $k$‑th iterate appears in the snapshot set, and there are at
    most $k$ iterates between two consecutive snapshots.)

\item \textbf{Combining (7)–(11).}  
    Multiply (7) by $\frac{2}{\eta T}$ and use (8)–(11):
    \[
    \frac{1}{T}\sum_{t=0}^{T-1}\mathbb{E}\big[\|\nabla f(w_t)\|^{2}\big]
    \le
    \frac{2\big(f(w_0)-f^{\star}\big)}{\eta T}
    +L\eta G^{2}.
    \tag{12}
    \]
    Here $f^{\star}\le\mathbb{E}[f(w_T)]$.
\item \textbf{From iterate average to interpolated point.}  
    The interpolated weight $\tilde w$ defined in (3) is a convex
    combination of two consecutive snapshots $w^{(s_{-})}$ and $w^{(s_{+})}$,
    which are at most $k$ steps apart.  By $L$‑smoothness,
    \[
    \|\nabla f(\tilde w)-\nabla f(w^{(s_{-})})\|
    \le L\|\tilde w-w^{(s_{-})}\|
    \le L\frac{k\eta G}{\;}\quad\text{(using a worst‑case bound on one SGD step)}.
    \tag{13}
    \]
    Squaring and taking expectations yields
    \[
    \mathbb{E}\big[\|\nabla f(\tilde w)\|^{2}\big]
    \le 2\mathbb{E}\big[\|\nabla f(w^{(s_{-})})\|^{2}\big]
        +2L^{2}\eta^{2}k^{2}G^{2}.
    \tag{14}
    \]
\item \textbf{Bounding the snapshot gradient term.}  
    Using (10)–(12) with $S\approx T/k$,
    \[
    \mathbb{E}\big[\|\nabla f(w^{(s_{-})})\|^{2}\big]
    \le \frac{k}{T}\sum_{t=0}^{T-1}\mathbb{E}\big[\|\nabla f(w_t)\|^{2}\big]
    \le
    \frac{2k\big(f(w_0)-f^{\star}\big)}{\eta T}
    +kL\eta G^{2}.
    \tag{15}
    \]
\item \textbf{Final bound.}  
    Substitute (15) into (14) and simplify:
    \[
    \mathbb{E}\big[\|\nabla f(\tilde w)\|^{2}\big]
    \le
    \frac{4k\big(f(w_0)-f^{\star}\big)}{\eta T}
    +2kL\eta G^{2}
    +2L^{2}\eta^{2}k^{2}G^{2}.
    \tag{16}
    \]
    Because $k$ is a fixed constant, we absorb factors of $k$ into the
    constants and obtain the cleaner statement (4):
    \[
    \mathbb{E}\big[\|\nabla f(\tilde w)\|^{2}\big]
    \le
    \frac{2\big(f(w_0)-f^{\star}\big)}{\eta T}
    +L\eta G^{2}
    +\frac{4L^{2}\eta^{2}\sigma^{2}}{k}.
    \qedhere
    \tag{17}
    \]
\end{enumerate}

%=====================================================================
\section*{Rate Derivation (With Constants)}
Choose $\eta = \frac{c}{\sqrt{T}}$ with $c\le \frac{1}{L}$.
Plugging into (4) gives  
\[
\mathbb{E}\big[\|\nabla f(\tilde w)\|^{2}\big]
\le
\frac{2\big(f(w_0)-f^{\star}\big)}{c}\frac{1}{\sqrt{T}}
+ \frac{LcG^{2}}{\sqrt{T}}
+ \frac{4L^{2}c^{2}\sigma^{2}}{k}\frac{1}{T}.
\]
The dominant terms scale as $\mathcal{O}(T^{-1/2})$, establishing the
claimed rate.

%=====================================================================
\section*{Special Cases}
\begin{enumerate}[label=(\alph*)]
\item \textbf{Convex $f$.}  
    If $f$ is convex, the same derivation yields a bound on the
    suboptimality $ \mathbb{E}[f(\tilde w)-f^{\star}]$ of order
    $\mathcal{O}(1/\sqrt{T})$, matching classical SGD results.
\item \textbf{Deterministic gradients ($\sigma=0$).}  
    The third term in (4) disappears, and with $\eta=1/L$ we obtain
    $\mathbb{E}[\|\nabla f(\tilde w)\|^{2}]\le
    \frac{2L\big(f(w_0)-f^{\star}\big)}{T}$, i.e. an $\mathcal{O}(1/T)$ rate.
\item \textbf{Large snapshot period $k\to\infty$.}  
    As $k$ grows, the interpolation bias term $\frac{4L^{2}\eta^{2}\sigma^{2}}{k}$
    vanishes, and ISA behaves like ordinary SGD with a final iterate
    guarantee.
\end{enumerate}

%=====================================================================
\end{document}